\documentclass[11pt,a4paper]{article}

% Pakete
\usepackage[utf8]{inputenc}
\usepackage[T1]{fontenc}
\usepackage[ngerman]{babel}
\usepackage[left=18mm,right=18mm,top=18mm,bottom=18mm]{geometry}
\usepackage{helvet}
\renewcommand{\familydefault}{\sfdefault}
\usepackage{amsmath,amssymb}
\usepackage{graphicx}
\usepackage{tikz}
\usetikzlibrary{arrows.meta,positioning}
\usepackage{tcolorbox}
\tcbuselibrary{skins,breakable}
\usepackage{enumitem}
\usepackage{tabularx}
\usepackage{colortbl}
\usepackage{xcolor}
\usepackage[hidelinks]{hyperref}

% Fachfarbe Physik
\definecolor{physik}{HTML}{2c5aa0}
\definecolor{physiklight}{HTML}{e8f0fa}
\definecolor{loesung}{HTML}{c62828}

% Box-Stile
\tcbset{
    formelbox/.style={
        colback=white,
        colframe=physik,
        boxrule=2pt,
        arc=2pt,
        left=8pt,
        right=8pt,
        top=8pt,
        bottom=8pt
    },
    loesungbox/.style={
        colback=red!5,
        colframe=loesung,
        boxrule=1pt,
        arc=2pt,
        left=6pt,
        right=6pt,
        top=6pt,
        bottom=6pt
    }
}

% Kopfzeile
\usepackage{fancyhdr}
\pagestyle{fancy}
\fancyhf{}
\lhead{\small Physik 12 GK}
\chead{\small \textcolor{loesung}{\textbf{MUSTERLÖSUNG}}}
\rhead{\small Bohr-Atommodell}
\renewcommand{\headrulewidth}{0.4pt}

% Keine Einrückung
\setlength{\parindent}{0pt}
\setlength{\parskip}{6pt}

\begin{document}

% Titel
\begin{minipage}[t]{0.65\textwidth}
\vspace{0pt}
{\Large\textbf{Bohr-Atommodell}} \hfill \textcolor{loesung}{\textbf{MUSTERLÖSUNG}}\\[0.3em]
{\normalsize Mi 28.01.2026 | Physik 12 GK}
\end{minipage}

\vspace{0.3em}
\hrule height 1pt
\vspace{0.6em}

% ============================================================
% AUFGABE 1: Das Problem
% ============================================================

\textbf{Aufgabe 1: Das Problem des Rutherford-Modells} \hfill (3 BE)

\vspace{0.3em}

a) Beobachte die Simulation im Lernpfad. Beschreibe, was mit dem Elektron passiert und warum. (2 BE)

\begin{tcolorbox}[loesungbox]
Das Elektron auf der Kreisbahn ist eine \textbf{beschleunigte Ladung}. Nach der klassischen Elektrodynamik (Maxwell) muss es daher \textbf{elektromagnetische Strahlung} aussenden. Dabei verliert es kontinuierlich Energie und \textbf{spiralt in den Atomkern}.
\end{tcolorbox}

b) In welcher Größenordnung liegt die vorhergesagte Zeit bis zum Kollaps? (1 BE)

\begin{tcolorbox}[loesungbox]
$t \approx \mathbf{10^{-11} \text{ s}}$ (etwa 10 Pikosekunden)
\end{tcolorbox}

\vspace{0.5em}

% ============================================================
% AUFGABE 2: Spektrum
% ============================================================

\textbf{Aufgabe 2: Wasserstoff-Spektrum} \hfill (3 BE)

\vspace{0.3em}

a) Zeichne das Wasserstoff-Spektrum (Balmer-Serie) in den Balken ein:

\vspace{0.3em}

\begin{center}
\begin{tikzpicture}
    \draw[black, thick] (0,0) rectangle (12,0.8);
    % Lösungs-Linien
    \fill[violet!90] (0.34,0.05) rectangle (0.44,0.75);  % 410 nm
    \fill[violet!70] (1.13,0.05) rectangle (1.23,0.75);  % 434 nm
    \fill[cyan!80] (2.87,0.05) rectangle (2.97,0.75);    % 486 nm
    \fill[red!90] (8.54,0.05) rectangle (8.74,0.75);     % 656 nm
    % Beschriftung
    \node[below] at (0,0) {\scriptsize 400 nm};
    \node[below] at (4,0) {\scriptsize 500 nm};
    \node[below] at (8,0) {\scriptsize 600 nm};
    \node[below] at (12,0) {\scriptsize 700 nm};
    % Lösungs-Labels
    \node[above, font=\tiny\color{loesung}] at (0.39,0.8) {410};
    \node[above, font=\tiny\color{loesung}] at (1.18,0.8) {434};
    \node[above, font=\tiny\color{loesung}] at (2.92,0.8) {486};
    \node[above, font=\tiny\color{loesung}] at (8.64,0.8) {656};
\end{tikzpicture}
\end{center}

\vspace{0.3em}

b) Das Spektrum ist \hspace{0.3cm} $\square$ kontinuierlich \hspace{0.5cm} \textcolor{loesung}{$\blacksquare$} \textbf{diskret}

\vspace{0.3em}

c) Erkläre, was dieses Spektrum über die Energie im Atom aussagt. (2 BE)

\begin{tcolorbox}[loesungbox]
Das diskrete Spektrum zeigt, dass die Energie im Atom \textbf{nur bestimmte Werte} annehmen kann. Die Energie ist \textbf{quantisiert}. Bei Übergängen zwischen Energieniveaus werden nur Photonen mit bestimmten Energien (= bestimmten Wellenlängen) emittiert.
\end{tcolorbox}

\vspace{0.5em}

% ============================================================
% FORMELKASTEN
% ============================================================

\begin{tcolorbox}[formelbox]
\textbf{Energieniveaus im Wasserstoffatom (Bohr-Modell)}

\vspace{0.5em}

\begin{center}
{\Large $E_n = \dfrac{-13{,}6 \text{ eV}}{n^2}$} \hspace{2cm} $n = 1, 2, 3, ...$
\end{center}

\vspace{0.3em}

\textbf{Photonenenergie bei Übergang:} \hspace{1cm} $E_{\text{Photon}} = |E_{n_1} - E_{n_2}| = h \cdot f$

\vspace{0.3em}

\textbf{Wellenlänge:} \hspace{1cm} $\lambda = \dfrac{h \cdot c}{E_{\text{Photon}}}$ \hspace{1cm} mit $h \cdot c = 1240 \text{ eV} \cdot \text{nm}$
\end{tcolorbox}

\vspace{0.5em}

% ============================================================
% AUFGABE 3: Berechnungen
% ============================================================

\textbf{Aufgabe 3: Energieniveaus berechnen} \hfill (4 BE)

\vspace{0.3em}

a) Berechne die Energien $E_1$, $E_2$ und $E_3$. Trage die Werte in die Tabelle ein. (2 BE)

\vspace{0.3em}

\begin{center}
\begin{tabular}{|c|c|c|c|c|c|}
\hline
$n$ & 1 & 2 & 3 & 4 & $\infty$ \\
\hline
$E_n$ (eV) & \textcolor{loesung}{\textbf{-13,6}} & \textcolor{loesung}{\textbf{-3,4}} & \textcolor{loesung}{\textbf{-1,51}} & $-0{,}85$ & 0 \\
\hline
\end{tabular}
\end{center}

\vspace{0.3em}

\begin{tcolorbox}[loesungbox]
\textbf{Rechnung:}
\begin{align*}
E_1 &= \frac{-13{,}6 \text{ eV}}{1^2} = \mathbf{-13{,}6 \text{ eV}} \\[0.3em]
E_2 &= \frac{-13{,}6 \text{ eV}}{2^2} = \frac{-13{,}6 \text{ eV}}{4} = \mathbf{-3{,}4 \text{ eV}} \\[0.3em]
E_3 &= \frac{-13{,}6 \text{ eV}}{3^2} = \frac{-13{,}6 \text{ eV}}{9} = \mathbf{-1{,}51 \text{ eV}}
\end{align*}
\end{tcolorbox}

\vspace{0.5em}

b) Berechne die Energie des Photons beim Übergang $n = 3 \to n = 2$. (2 BE)

\begin{tcolorbox}[loesungbox]
\textbf{Rechnung:}
\begin{align*}
E_{\text{Photon}} &= |E_3 - E_2| = |(-1{,}51 \text{ eV}) - (-3{,}4 \text{ eV})| \\[0.3em]
&= |-1{,}51 + 3{,}4| \text{ eV} = \mathbf{1{,}89 \text{ eV}}
\end{align*}
\end{tcolorbox}

\vspace{0.5em}

% ============================================================
% AUFGABE 4: Anwendung
% ============================================================

\textbf{Aufgabe 4: Wellenlänge berechnen} \hfill (3 BE)

\vspace{0.3em}

Berechne die Wellenlänge des Photons aus Aufgabe 3b. In welchem Farbbereich liegt diese Linie?

\begin{tcolorbox}[loesungbox]
\textbf{Rechnung:}
\begin{align*}
\lambda &= \frac{h \cdot c}{E_{\text{Photon}}} = \frac{1240 \text{ eV} \cdot \text{nm}}{1{,}89 \text{ eV}} = \mathbf{656 \text{ nm}}
\end{align*}

\textbf{Farbbereich:} \textbf{Rot} (sichtbares Licht, 620-700 nm)

Dies ist die H$_\alpha$-Linie der Balmer-Serie.
\end{tcolorbox}

\vspace{0.5em}

% ============================================================
% AUFGABE 5: Transfer
% ============================================================

\textbf{Aufgabe 5: Ionisierungsenergie} \hfill (2 BE)

\vspace{0.3em}

Berechne die Energie, die nötig ist, um das Elektron aus dem Grundzustand ($n=1$) vollständig zu entfernen.

\begin{tcolorbox}[loesungbox]
\textbf{Rechnung:}
\begin{align*}
E_{\text{Ion}} &= |E_{\infty} - E_1| = |0 - (-13{,}6 \text{ eV})| = \mathbf{13{,}6 \text{ eV}}
\end{align*}
\end{tcolorbox}

\vspace{0.5em}

% ============================================================
% AUFGABE 6: Transfer (AFB III)
% ============================================================

\textbf{Aufgabe 6: Transfer -- Helium-Ion} \hfill (3 BE)

\vspace{0.3em}

Das einfach ionisierte Helium-Atom (He$^+$) enthält nur ein Elektron, aber der Kern hat die Ladung $Z = 2$. Die Energieformel lautet dann:

\begin{center}
$E_n = \dfrac{-13{,}6 \text{ eV} \cdot Z^2}{n^2}$
\end{center}

Berechne die Wellenlänge des Photons beim Übergang $n = 3 \to n = 2$ in He$^+$.

\begin{tcolorbox}[loesungbox]
\textbf{Rechnung:}
\begin{align*}
E_3(\text{He}^+) &= \frac{-13{,}6 \text{ eV} \cdot 2^2}{3^2} = \frac{-13{,}6 \cdot 4}{9} \text{ eV} = \mathbf{-6{,}04 \text{ eV}} \\[0.3em]
E_2(\text{He}^+) &= \frac{-13{,}6 \text{ eV} \cdot 4}{4} = \mathbf{-13{,}6 \text{ eV}} \\[0.3em]
E_{\text{Photon}} &= |E_3 - E_2| = |-6{,}04 - (-13{,}6)| \text{ eV} = \mathbf{7{,}56 \text{ eV}} \\[0.3em]
\lambda &= \frac{1240 \text{ eV} \cdot \text{nm}}{7{,}56 \text{ eV}} = \mathbf{164 \text{ nm}} \quad \text{(UV-Strahlung)}
\end{align*}
\end{tcolorbox}

\vspace{0.5em}

% ============================================================
% AUFGABE 7: Beurteilung (AFB III)
% ============================================================

\textbf{Aufgabe 7: Grenzen des Bohr-Modells} \hfill (3 BE)

\vspace{0.3em}

a) Das Bohr-Modell kann die Spektren von Atomen mit mehr als einem Elektron (z.\,B. Helium) nicht korrekt vorhersagen. Erkläre, warum das Modell hier an seine Grenzen stößt. (2 BE)

\begin{tcolorbox}[loesungbox]
Das Bohr-Modell berücksichtigt nur die \textbf{Coulomb-Wechselwirkung} zwischen dem Kern und \emph{einem} Elektron. Bei Mehrelektronenatomen gibt es zusätzlich die \textbf{Elektron-Elektron-Abstoßung}, die das Modell nicht erfasst. Außerdem sind die Bahnen im quantenmechanischen Bild keine scharfen Kreise, sondern \textbf{Aufenthaltswahrscheinlichkeiten} (Orbitale).
\end{tcolorbox}

b) Trotz dieser Einschränkung wird das Bohr-Modell noch heute im Unterricht verwendet. Beurteile, ob dies sinnvoll ist. (1 BE)

\begin{tcolorbox}[loesungbox]
\textbf{Ja, sinnvoll:}
\begin{itemize}[noitemsep,topsep=2pt]
\item Das Modell erklärt anschaulich das Konzept der \textbf{Quantisierung} und diskreter Spektren.
\item Es liefert korrekte Ergebnisse für wasserstoffähnliche Atome (H, He$^+$, Li$^{2+}$).
\item Es fördert \textbf{Modellkompetenz}: Schüler lernen, dass Modelle begrenzt sind und durch bessere ersetzt werden.
\end{itemize}
\end{tcolorbox}

\vfill

\begin{center}
\rule{10cm}{0.5pt}

\textbf{Gesamt:} \textcolor{loesung}{\textbf{21}} / 21 BE
\end{center}

\end{document}
